\subsection{Метод Марквардта}
\label{subsec:marquardt}

\textbf{Идея:} соединить характеристики методов первого и второго порядков.

\begin{idea}
В литературе этот метод чаще называют \textbf{методом
Левенберга--Марквардта}, а замену $H_k$ на $H_k + \lambda_k E$ —
\textbf{регуляризацией} гессиана. Это третья (наряду с двумя из
\S\ref{subsec:newton-variants}) стандартная модификация метода Ньютона, и
единственная из них, которая чинит не стоимость итерации, а саму матрицу: она
применима и тогда, когда $H_k$ не положительно определена.
\end{idea}

\begin{algo}[title={Метод Марквардта}]
\textbf{Начальный этап:} задать $\eps > 0$, $\lambda_0 = 10^4$, начальное
приближение $x_0$, $k = 0$.

\medskip
\textbf{Основной этап:}
\begin{steps}
  \item Вычислить $\nabla f(x_k)$.
  \item Проверить $\norm{\nabla f(x_k)} < \eps$: \ да — конец; нет — к~шагу~3.
  \item Вычислить $H_k$ и найти $p_k$ из системы
        \[
          \big(H_k + \lambda_k E\big)\, p_k = -\nabla f(x_k) .
        \]
  \item $x_{k+1} := x_k + p_k$.
  \item Проверить $f(x_{k+1}) < f(x_k)$: \ да — к~шагу~6; нет — к~шагу~7.
  \item $\lambda_{k+1} := \dfrac{\lambda_k}{2}$, \ $k := k+1$, к~шагу~1.
  \item $\lambda_{k+1} := 2\lambda_k$, вернуться к~шагу~3.
\end{steps}
\end{algo}

\begin{idea}
Смысл параметра $\lambda_k$ виден из крайних случаев системы на шаге~3:
\begin{itemize}[nosep,left=0pt]
  \item $\lambda_k \to 0$: система превращается в $H_k p_k = -\nabla f(x_k)$ —
        чистый метод Ньютона, быстрый вблизи минимума;
  \item $\lambda_k \to \infty$: слагаемое $\lambda_k E$ доминирует, и
        $p_k \approx -\frac{1}{\lambda_k}\nabla f(x_k)$ — градиентный шаг с
        малой длиной, надёжный вдали от минимума.
\end{itemize}
Старт с $\lambda_0 = 10^4$ — это «начать как градиентный метод». Удачный шаг
уменьшает $\lambda$ вдвое (сдвиг в сторону Ньютона), неудачный — удваивает
(откат к осторожному градиентному шагу). Попутно $\lambda_k E$ делает матрицу
положительно определённой даже там, где $H_k$ таковой не является, то есть
направление $p_k$ всегда остаётся направлением спуска.

В шаге~7 возврат именно к~шагу~3, а не к~4: после изменения $\lambda_k$ систему
надо решать заново, иначе $p_k$ останется прежним и точка не сдвинется.
\end{idea}

\subsection{Методы сопряжённых градиентов}
\label{subsec:conjugate-gradients}

\[
  x^{(k+1)} = x^{(k)} - \alpha^{(k)} p^{(k)} , \qquad k \in \NN_0 ,
\]
\[
  p^{(0)} = \nabla f\big(x^{(0)}\big) ,
  \qquad
  p^{(k)} = \nabla f\big(x^{(k)}\big) + \beta^{(k)} p^{(k-1)} ,
  \qquad k \in \NN ,
\]
\[
  \beta^{(k)}
  = \frac{\norm{\nabla f\big(x^{(k)}\big)}^2}
         {\norm{\nabla f\big(x^{(k-1)}\big)}^2}
  = \left( \sum_{i=1}^{n}
      \left(\frac{\partial f\big(x^{(k-1)}\big)}{\partial x_i}\right)^{\!2}
    \right)^{\!-1}
    \sum_{i=1}^{n}
      \left(\frac{\partial f\big(x^{(k)}\big)}{\partial x_i}\right)^{\!2} ,
  \qquad k \in \NN ,
\]
\[
  \beta^{(mn)} = 0 , \qquad m \in \NN_0 .
\]

\begin{idea}[title={Что означает $\beta^{(mn)} = 0$}]
Это \textbf{рестарт}: каждые $n$ итераций коэффициент обнуляют, и очередное
направление берётся чисто градиентным, $p^{(mn)} = \nabla f\big(x^{(mn)}\big)$.

Причина та же, что у сброса $A_k = E$ в квазиньютоновских методах. Для
квадратичной функции $n$ направлений $p^{(0)}, \dots, p^{(n-1)}$ взаимно
сопряжены относительно гессиана, и метод приходит в точный минимум не более
чем за $n$ шагов — дальше строить сопряжённые направления просто не из чего.
Для неквадратичной $f$ гессиан меняется от точки к точке, накопленная
сопряжённость теряется, и продолжение рекуррентной формулы даёт направления,
всё хуже связанные с текущей кривизной. Рестарт каждые $n$ шагов обрывает
накопление ошибки и гарантирует, что метод не застрянет: направление спуска
заведомо корректно.

Записанная формула — выбор $\beta$ по \textbf{Флетчеру--Ривсу}; вторая строка
с суммами — та же дробь, расписанная покоординатно.
\end{idea}

\begin{idea}
Место метода среди прочих: производные второго порядка не нужны, хранить
матрицу $n \times n$ (как в квазиньютоновских методах) тоже не нужно — только
предыдущее направление. Поэтому метод сопряжённых градиентов — основной выбор
при очень больших $n$, где $H$ или $A_k$ просто не помещаются в память.
\end{idea}
